\documentclass[11pt]{article}
\usepackage[margin=1in]{geometry}
\usepackage{amsmath,amssymb,amsthm,hyperref}
\newtheorem{theorem}{Theorem}
\title{A bounded strictly starlike domain with a nonstarlike Cassinian ball}
\author{Principia Math in Progress}
\date{12 September 2026}
\begin{document}
\maketitle
\section*{Task 1. Resolve MathDB 331705}
For a proper domain $G\subset\mathbb R^n$, its Cassinian metric is
\[
 c_G(x,y)=\sup_{p\in\partial G}\frac{|x-y|}{|x-p||p-y|},
 \qquad B_{c_G}(x,r)=\{y\in G:c_G(x,y)<r\}.
\]
Conjecture 4.7 of \cite{KMS} asserts that if $G$ is starlike with respect to
$x$, then $B_{c_G}(x,r)$ is starlike with respect to $x$ for every $r>0$.
The following counterexample also satisfies the stricter boundary-segment
definition of starlikeness used in \cite{Klen}.

\begin{theorem}
In $\mathbb R^2$, set
\[
 K=\{(u,v):u\ge1,\ 4|v|\le u-1\},\qquad
 G=B(0,10)\setminus K.
\]
Then $G$ is a bounded domain, strictly starlike with respect to $0$, but
$B_{c_G}(0,5/4)$ is not starlike with respect to $0$.
\end{theorem}
\begin{proof}
The set $K$ is closed, so $G$ is open, bounded, and contains $0$.
For any $w=(u,v)\in\overline G$ and $0\le\tau<1$, we have
$|\tau w|<10$. If $\tau u<1$, then $\tau w\notin K$.
Otherwise $u\ge1$ and $w\notin\operatorname{int}K$ imply
$4|v|\ge u-1$, whence
\[
 4|\tau v|\ge\tau(u-1)=\tau u-\tau>\tau u-1.
\]
Again $\tau w\notin K$. Thus $[0,w)\subset G$ for all $w\in\overline G$;
in particular $G$ is path connected and strictly starlike about $0$.

Take $y=(4,4)$ and $z=(1,1)=y/4$. Both points lie in $G$.
We compute their exact Cassinian distances from $0$.
Every boundary point belongs either to the circle $|p|=10$ or to one of
the truncated rays
\[
 p_\pm(t)=(1+t,\,\pm t/4),\qquad t\ge0.
\]
For $w=(a,a)$ with $a\ge0$, the squared distances satisfy
\[
 |p_-(t)-w|^2-|p_+(t)-w|^2=at\ge0,
 \qquad |p_-(t)|=|p_+(t)|.
\]
It suffices to bound the denominator on the upper ray. Put
\[
 P_a(t)=|p_+(t)|^2|p_+(t)-(a,a)|^2
 =\left((1+t)^2+\frac{t^2}{16}\right)
  \left((1+t-a)^2+\left(\frac t4-a\right)^2\right).
\]
Direct expansion and completion of the square give
\begin{align*}
 P_4(t)-25&=42t+t^2\left[
    \frac{289}{256}\left(t-\frac{48}{17}\right)^2+\frac{21}{8}\right]\ge0,\\
 P_1(t)-1&=\frac32t+\frac98t^2+\frac{51}{32}t^3+
                \frac{289}{256}t^4\ge0.
\end{align*}
Equality holds at $t=0$, the actual boundary point $p=(1,0)$.
Therefore the supremum over the ray portions is $4\sqrt2/5$ for $y$,
and $\sqrt2$ for $z$.

The circular portion cannot increase either value. Indeed, for $|p|=10$,
\[
 \frac{|w|}{|p||p-w|}\le\frac{|w|}{10(10-|w|)}.
\]
Since $|y|=4\sqrt2<6$ and $|z|=\sqrt2<2$, this upper bound is respectively
less than $3/20<4\sqrt2/5$ and $1/40<\sqrt2$. Consequently
\[
 c_G(0,y)=\frac{4\sqrt2}{5}<\frac54,
 \qquad c_G(0,z)=\sqrt2>\frac54.
\]
The squared comparisons are $32/25<25/16<2$.
Thus $y\in B_{c_G}(0,5/4)$ while the point $z=y/4$ of the segment
$[0,y]$ is outside that ball, proving the assertion.
\end{proof}

\paragraph{Scope.}
This refutes the universal claim in dimension two, even with a bounded
strictly starlike domain whose removed set has positive interior. It does
not resolve the separate convex-domain conjecture adjacent to the selected
statement in \cite{KMS}.

\paragraph{Task 2. Smooth the domain.}
Construct a counterexample with a smooth strictly starlike boundary,
using quantitative control of the two strict Cassinian inequalities.

\paragraph{Task 3. Determine the sharp universal small-radius guarantee.}
Determine the largest universal constant $a$, if one exists, such that
$r\,\operatorname{dist}(x,\partial G)\le a$ guarantees starlikeness of
$B_{c_G}(x,r)$ for every starlike domain $G$ centered at $x$.

\paragraph{Task 4. Characterize centers with radial monotonicity.}
Characterize geometric conditions on $(G,x)$ that ensure
$t\mapsto c_G(x,x+tv)$ is nondecreasing along every segment in $G$.

\paragraph{Task 5. Classify the cone family.}
For $G=B(0,R)\setminus\{(u,v):u\ge1,\ |v|\le b(u-1)\}$,
determine the parameter and radius ranges producing nonstarlike
Cassinian balls about $0$.

\begin{thebibliography}{9}
\bibitem{KMS} R. Kl\'en, M. R. Mohapatra and S. K. Sahoo,
\emph{Geometric properties of the Cassinian metric},
\href{https://arxiv.org/abs/1511.01298}{arXiv:1511.01298v1} (2015),
Definition on page 2 and Conjecture 4.7 on page 8.
Published in Math. Nachr. 290 (2017), 1531--1543,
\href{https://doi.org/10.1002/mana.201600117}{doi:10.1002/mana.201600117}.
\bibitem{Klen} R. Kl\'en, \emph{Hyperbolic type distances in starlike domains},
\href{https://arxiv.org/abs/1603.00590}{arXiv:1603.00590v1} (2016), page 3.
\end{thebibliography}
\end{document}
